\chapter{Conclusions and Future Work}
\label{ch:conclusions}

This thesis assessed whether building attributes inferred from \ac{SVI} can substitute for reference attributes in a TABULA-based building-stock enrichment workflow.
The traceable semi-synthetic dataset enabled attribute extraction, archetype assignment, physically weighted archetype differences, and downstream energy-class classification to be evaluated within one controlled framework.
The findings provide bounded answers to the three research objectives.

\section{Conclusions}
\label{sec:conclusions}

For Objective~1, \ac{SVI} provided predictive information for building type, construction year, and floor count, although the quality of extraction differed across attributes and classes.
Construction year showed the strongest building-level performance.
Building-type accuracy was influenced by the dominant Apartment Block class, and the floor-count result was bounded by the narrow range represented in the sample.
Among the evaluated configurations, the frozen DINOv2 backbone with supervised prediction layers provided the most consistent overall results while requiring substantially fewer parameter updates than full ResNet-50 fine-tuning.
The zero-shot InternVL3-2B configuration did not provide sufficient accuracy for the three-attribute extraction task under the evaluated protocol.

For Objective~2, \ac{SVI}-derived attributes reproduced the dominant TABULA-NL archetype cells but did not provide balanced coverage of the complete matrix.
The best joint-cell accuracy exceeded the majority-cell baseline by only approximately 3.5 percentage points, and recall remained low for sparsely represented cells.
Physically weighted archetype-difference scoring distinguished the trained configurations more clearly from the zero-shot configuration, but the trained configurations still overestimated stock-level transmission heat loss by approximately 10\%.
These results also separate two uses of the archetype assignment.
The tabulated U-values add no independent information to energy-class classification once building type and construction period are present, whereas heat-loss estimation depends directly on the accuracy and bias of the assigned physical parameters.
The comparison remains relative to reference archetypes and does not constitute validation against measured building performance.

For Objective~3, \ac{SVI}-derived attributes provided a conditional partial substitute when reference attributes were unavailable, but they were not an equivalent replacement.
Under the pooled protocol with fixed buildings, data splits, and downstream classifier, replacing the reference attributes with DINOv2-derived attributes reduced binary macro-F1 by 0.002.
However, both routes achieved D--G recall of approximately 0.11, so the small difference between them does not imply strong binary classification.
The substitution cost was more apparent in the seven-class diagnostic task, where macro-F1 decreased by 0.022.
Feature ablation identified construction year as the largest marginal contributor among the existing downstream attributes, making construction-year extraction the first priority for improving the attribute-mediated route.

The Amsterdam-held-out evaluation further limits the scope of this conclusion.
Excluding Amsterdam from training increased construction-year error, while both the reference and attribute-mediated routes approached the uniform-random expectation on the concentrated Amsterdam target sample.
The pooled result therefore cannot be extended directly to a city excluded from training.
Cross-city use requires separate evaluation of attribute transfer and target-sample composition.

Overall, \ac{SVI}-derived attributes can supplement missing reference metadata in a TABULA-based building-stock workflow, but their suitability depends on the attribute, downstream task, and relationship between the training and target populations.
Within the four-city sample evaluated here, frozen DINOv2 offered the most consistent source of the three visual attributes, and construction year provided the clearest downstream utility.
The controlled comparison establishes how attribute contribution, substitution cost, physically weighted archetype difference, and cross-city performance can be evaluated separately rather than treating \ac{SVI}-derived information as either universally usable or unusable.

\section{Contributions}
\label{sec:contributions}

This thesis makes three contributions to the evaluation of \ac{SVI}-derived attributes for building-stock enrichment.
These contributions concern the controlled assessment of attribute sources and their downstream consequences rather than the development of a new predictive architecture.

\textbf{Controlled evaluation of SVI-based attribute extraction.}
A common building population and attribute schema are used to compare a fully fine-tuned ResNet-50, a frozen DINOv2 backbone with supervised prediction layers, and zero-shot InternVL3-2B.
Evaluating all configurations on the same hold-out buildings separates performance differences among construction year, building type, and floor count and identifies the relative strengths and limitations of the three inference strategies.

\textbf{Traceable propagation from attributes to archetype-based physical estimates.}
The building-level dataset links Mapillary imagery, reference attributes, reconstructed 3DBAG geometry, EPC records, and TABULA-NL archetypes while retaining the provenance of each value.
Predicted attributes are propagated through joint archetype assignment to the transmission heat-transfer coefficient $H'_{tr}$.
This structure distinguishes extraction errors from their consequences for archetype coverage, physical parameter assignment, and stock-level heat-loss bias.

\textbf{Controlled assessment of downstream substitution and its applicability.}
Reference and \ac{SVI}-derived attributes are compared using the same buildings, data splits, and downstream classifier.
The resulting comparison attributes the observed performance change to the replacement of the attribute source rather than to a different evaluation population or classifier.
Feature ablation identifies the relative downstream contribution of the existing attributes, while the sample-composition and Amsterdam-held-out evaluations define the conditions under which the pooled substitution result does and does not apply.

\section{Limitations}
\label{sec:limitations}

\textbf{Data coverage, energy-class targets, and sample composition.}
Mapillary is a crowdsourced source with uneven spatial coverage, acquisition dates, viewpoints, and image quality.
Buildings without suitable street-level views are less likely to enter the experimental sample.
The resulting sample is dominated by Amsterdam and apartment blocks, while several construction periods, building types, archetype cells, and energy classes remain sparsely represented.
The diagnostic analyses show that this composition affects class-balanced performance and cannot be addressed by increasing sample count while retaining the same distribution.
Energy-class targets introduce further uncertainty because one building can be linked to multiple certificates that differ across units or registration dates.
The reported agreement statistics describe the adopted aggregation rules and are not strict upper bounds on prediction.
Mapillary acquisition dates were not systematically aligned with EPC registration dates, so the imagery and certificate may represent different building states.

\textbf{Scope of model development.}
The study trained supervised models and prediction layers to compare predefined inference configurations, but its purpose was not to introduce a new architecture or maximise attainable prediction performance.
It did not exhaustively evaluate model size, pre-training strategy, hyperparameters, training duration, or alternative image-selection procedures.
InternVL3-2B was evaluated only without parameter updates, and no \ac{VLM} was fine-tuned on the study dataset.
The reported accuracy and macro-F1 values therefore characterise the selected configurations rather than an upper bound on what can be inferred from \ac{SVI}.
Additional representative training data, task-specific fine-tuning, or broader optimisation may improve performance, but the present experiments do not determine the magnitude of such gains.
The sample-composition analysis also shows that adding data without improving representation is not sufficient by itself.

\textbf{Archetype and physical assumptions.}
Converting reference or predicted attributes into TABULA-NL archetypes requires operational mapping decisions, including the assignment of semi-detached houses and flats to broader building types.
These decisions were not examined through separate sensitivity tests.
\Ac{SVI} can show visible facade or window changes but cannot reliably reveal concealed insulation, heating-system upgrades, or unit-level renovation.
The assigned TABULA parameters therefore represent archetypal as-built characteristics rather than verified current envelope conditions.
In addition, the $H'_{tr}$ analysis compares \ac{SVI}-derived and reference-derived archetype assignments and is not a validation against measured heat loss or energy consumption.

\textbf{Geographical scope.}
The empirical evaluation covers four cities in the Netherlands, and Amsterdam is the only fully held-out target city.
The Amsterdam target sample is concentrated in apartment blocks and one dominant archetype cell, so its result is not an average estimate of transfer to other cities.
The use of national registries, energy-class definitions, 3DBAG geometry, and the TABULA-NL taxonomy also limits direct transfer to other countries.
Cross-country application requires country-specific assessment of data sources, archetype mappings, target definitions, and model performance.

\section{Future Work}
\label{sec:future_work}

\textbf{Improve data coverage and target alignment.}
Future datasets should increase representation of sparsely sampled construction periods, residential building types, archetype cells, energy classes, and urban contexts.
Expansion should be guided by these groups rather than preserve the current sample composition.
Image acquisition and EPC registration dates should be aligned where possible.
For apartment buildings, unit-level prediction and alternative building-level aggregation rules should be evaluated to address the mismatch between facade-level observation and unit-level certification.

\textbf{Optimise visual attribute extraction.}
Construction year is the first priority for further model development because it made the largest marginal contribution to downstream energy-class classification among the existing attributes.
A fine-tuned open-source \ac{VLM} should be compared with frozen DINOv2 and fully fine-tuned \ac{CNN} configurations using the same buildings, splits, attribute definitions, and downstream protocol.
Representative additional training data, class-balanced objectives, ordinal construction-period formulations, calibration, and systematic hyperparameter optimisation should also be evaluated.
These experiments would establish how much of the current error is associated with model capacity, training design, target quality, and the visual evidence available in \ac{SVI}.

\textbf{Extend geographical validation and adaptation.}
Leave-one-city-out evaluation should be repeated with each city as the target to distinguish general transfer behaviour from the composition of Amsterdam.
Separate experiments could then evaluate target-city calibration using a limited number of labelled buildings or suitable geospatial labels.
Cross-country evaluation should begin with explicit mappings among national registries, energy-class definitions, and TABULA taxonomies, followed by within-country validation before transfer between countries is assessed.

\textbf{Represent renovation and validate physical performance effects.}
Future workflows should use the source best suited to each attribute.
Reconstructed geospatial data should remain the preferred source for building geometry where equivalent 3D data are available, while \ac{SVI} analysis can focus on visible facade characteristics and potentially outdated or missing semantic attributes.
Renovation records, measured energy consumption, or building-specific thermal assessments are needed to distinguish as-built archetypes from current building conditions.
If the residual between measured consumption and archetype-based estimates is examined as a renovation signal, weather, occupancy, floor area, and heating systems must first be controlled.
These data would also enable the physical evaluation to move from comparison with reference archetypes to validation against observed building performance.
